\documentclass[12pt]{article}
\usepackage[margin=1in]{geometry}
\usepackage{amsmath,amssymb,amsthm}
\usepackage{graphicx}
\usepackage[hidelinks]{hyperref}
\hypersetup{pdftitle={Lorentzian Read-out from Local Unitary Propagation: Clifford Relations, Causal Cones and Spectral Mass},pdfauthor={Sidong Liu}}
\usepackage{booktabs}
\usepackage{array}
\newcolumntype{P}[1]{>{\raggedright\arraybackslash}p{#1}}

\newtheorem{theorem}{Theorem}
\newtheorem{lemma}[theorem]{Lemma}
\newtheorem{proposition}[theorem]{Proposition}
\newtheorem{corollary}[theorem]{Corollary}
\newtheorem{conjecture}[theorem]{Conjecture}
\theoremstyle{definition}
\newtheorem{definition}[theorem]{Definition}
\theoremstyle{remark}
\newtheorem{remark}[theorem]{Remark}

\newcommand{\tr}{\operatorname{tr}}
\newcommand{\conv}{\operatorname{conv}}
\newcommand{\Herm}{\operatorname{Herm}}

\title{Lorentzian Read-out from Local Unitary Propagation:\\
Clifford Relations, Causal Cones and Spectral Mass}

\author{Sidong Liu, PhD \\
\small iBioStratix Ltd \\
\small \texttt{sidongliu@hotmail.com}}

\date{October 2026}

\begin{document}
\emergencystretch=2em
\raggedbottom
\maketitle

% ============================================================
\begin{abstract}
A local unitary propagation law can supply the relative sign of a Lorentzian read-out without turning a positive control metric into an indefinite one. For a first-order Hermitian symbol with independent spatial coefficients, all bands share one quadratic shell exactly when those coefficients satisfy the corresponding Euclidean Clifford relations. The resulting form is $\omega^2-Q(q)$; it describes the full cone when both frequency branches occur. The evolution time is supplied. The internal dimension bounds the number of anticommuting directions, while calibrated coefficient errors give a forward certificate for approximate shells.

A shared shell for one species does not imply one cone for all matter. Fixed-block one-particle mixing cannot merge distinct shells. In a mediator model, sufficiently accumulated interactions align speeds, even when different edges are active on disjoint scale intervals. At a supplied Lorentz-invariant conformal fixed point, nonconserved spin-2 perturbations are irrelevant; additional conserved currents are the marginal exceptions. Neither argument creates the assumed infrared regime.

On a lattice, support, group and rest-and-curvature speeds differ. An actual record protocol certifies the angle of one coined walk and thereby its ideal dictionary parameters, not its noisy long-time transport. The interfaces to physical control and spectral mass require, respectively, calibrated generator alignment and a positive-energy representation. These conditions make the structural connection testable while leaving the formation of the carrier and accessible momentum window as dynamical questions.

\medskip
\noindent\textbf{Keywords}: Lorentzian signature, Clifford algebra, hyperbolic polynomials, quantum cellular automata, quantum walks, Weyl and Dirac equations, causal cone, spectral mass, Q-RAIF
\end{abstract}

% ============================================================
\section{Introduction}
\label{sec:intro}

Programmes that reconstruct geometry from quantum data usually take the Lorentzian signature as given. Holographic entanglement formulas and thermodynamic derivations of the Einstein equation relate entropies to areas and energy fluxes in a spacetime whose causal structure is already Lorentzian~\cite{RyuTakayanagi2006,Jacobson1995}. The Holographic Alaya-Field Framework (HAFF) studies which distinctions a physical reduction can resolve and what geometric data can be reconstructed from them~\cite{Liu2026PaperA,Liu2026PaperB}; a spacetime metric needs a further dictionary there as well. This paper asks where the indefinite direction of such a read-out can come from in a model whose dynamics is given explicitly.

Two facts fix the starting point. A positive definite form stays positive under every linear change of variables (Proposition~\ref{prop:sylvester}), so no control metric, information metric or work form becomes Lorentzian by reparameterization. The companion paper on persistence makes the same distinction physically: its descent metric and field-derived work form are positive. Under an additional all-state work condition and calibrated generator alignment it relates work to the positive spatial shell, not directly to the Lorentzian form~\cite{Liu2026QRAIF_B}. Second, imposing Clifford relations $\{e_\mu,e_\nu\}=2q_{\mu\nu}\mathbf 1$ for a chosen Lorentzian $q$ puts the signature into the premise; the universal property of $Cl(V,q)$ then supplies an encoding, rather than a signature-selection mechanism~\cite{Hestenes1966,Doran2003}.

The model class is local unitary propagation: a particle with an internal Hilbert space $\mathbb C^s$ hops on a lattice by a translation-invariant, finite-range unitary step, as in quantum walks and the one-particle sector of quantum cellular automata~\cite{BialynickiBirula1994,Meyer1996,Strauch2006,BisioDArianoTosini2015,DArianoPerinotti2014}. Its low-energy generator is a first-order Hermitian matrix symbol. In this setting four objects can be defined separately: a positive norm on generators, which plays the role of a control metric; the causal cones of the dynamics; the algebra of the symbol's coefficients; and the spectral mass. For a first-order Hermitian symbol, unitarity alone makes the characteristic polynomial hyperbolic in the frequency (Theorem~\ref{thm:hyperbolic}). For linearly independent spatial coefficients, all bands lie on one quadratic shell if and only if the coefficients satisfy the shell relations of Theorem~\ref{thm:shell}; the characteristic form is then $\omega^2$ minus the Hilbert--Schmidt norm of the spatial symbol and has signature $(1,d)$, and the Lorentzian Clifford generators are $\gamma^0=\beta$, $\gamma^i=\beta\alpha_i$ (Theorem~\ref{thm:shell}, Corollary~\ref{cor:gamma}). The internal dimension bounds the spatial one, $d+[\mu\neq0]\le2\nu_2(s)+1$. For a qubit carrier three is the largest dimension with an isotropic massless cone, and the only one in which every internal rotation axis is coupled to a direction of motion (Theorem~\ref{thm:dimension}, Proposition~\ref{prop:qubit}).

On the lattice, group velocities lie in the convex hull of the neighbourhood, so the speed of any massless read-out cone is at most the inradius of that polytope per step (Proposition~\ref{prop:polytope}). Exact formulas for the split-step walk and computations for the Weyl automaton then show that the read-out map is not fixed by the dynamics: a band point, a momentum window and the observable that defines the speed all have to be chosen (Section~\ref{sec:lattice}).

The mathematical ingredients are standard: Dirac's square root of the Klein--Gordon operator~\cite{Dirac1928}, G{\aa}rding's hyperbolic polynomials~\cite{Garding1959}, the representation theory of complex Clifford algebras~\cite{LawsonMichelsohn1989}, and the derivations of Weyl and Dirac equations from quantum walks and automata~\cite{BialynickiBirula1994,Strauch2006,DArianoPerinotti2014}. Hyperbolicity as the defining property of a physical dispersion relation is developed in Ref.~\cite{RaetzelRiveraSchuller2011}, and an information-theoretic argument linking three spatial dimensions to the qubit is given in Ref.~\cite{MullerMasanes2013}. This paper separates the four objects in one model, identifies which inputs of a Lorentzian read-out follow from propagation, and calculates how the remaining read-out choices enter. The several-species analysis adds a quantitative fixed-block mixing obstruction and applies cut-balanced convergence to interacting speeds that need not couple simultaneously.

The argument proceeds from the symbol and its shell to dimension, lattice windows and several-species dynamics. A residual bound connects the exact shell criterion to calibrated approximate symbols; a finite angle certificate then connects one such dictionary to actual records. The concluding sections distinguish these contributions from the physical selection of a carrier, an interacting fixed point and a complete source-and-return apparatus.

% ============================================================
\section{Local unitary propagation and four objects}
\label{sec:model}

\subsection{The model}

Let $\Lambda\subset\mathbb R^d$ be a lattice of physical positions and let the one-particle Hilbert space be $\ell^2(\Lambda)\otimes\mathbb C^s$. A translation-invariant step with finite neighbourhood $N\subset\Lambda$ acts as
\begin{equation}
(W\psi)(x)=\sum_{h\in N}A_h\,\psi(x-h),\qquad A_h\in M_s(\mathbb C),
\label{eq:step}
\end{equation}
and is unitary exactly when its symbol
\begin{equation}
W(k)=\sum_{h\in N}A_h\,e^{-ik\cdot h}
\label{eq:symbol}
\end{equation}
is unitary for every $k$ in the Brillouin zone. One step takes time $\tau$. The eigenvalues of $W(k)$ are $e^{-i\omega_b(k)\tau}$ with quasi-energies $\omega_b(k)$ defined modulo $2\pi/\tau$, and the group velocity of band $b$ is $\nabla_k\omega_b$. With this convention the pure shift $A_{h}=\mathbf 1$ for a single $h$ moves every state by $+h$ per step.

Near a point $k_*$ where a gap closes or nearly closes, the step is described by an effective generator. If a parameter $\theta$ controls the gap, we say that the walk has \emph{first-order symbol}
\begin{equation}
H(q)=\sum_{i=1}^d q_i\alpha_i+\mu\beta,\qquad \alpha_i,\beta\in\Herm(s),
\label{eq:firstorder}
\end{equation}
at $k_*$ if $(i/\tau)\operatorname{Log}W_\theta(k_*+q)=H(q)+O(|q|^2+|\mu|^2)$ with $\mu$ proportional to $\theta$. This is the scaling limit studied for quantum walks and automata~\cite{Strauch2006,BisioDArianoTosini2015,DArianoPerinotti2014}: $H$ generates the continuum dynamics $i\partial_t\psi=H(-i\nabla)\psi$ when lattice spacing and step time go to zero at fixed ratio. For example, the split-step walk of Section~\ref{sec:walk} has $W_\theta(k)=\operatorname{diag}(e^{-ika},e^{ika})R(\theta)$ with $R(\theta)=\cos\theta-i\sin\theta\,\sigma_y$, and at $k_*=0$ its first-order symbol is $H(q)=(a/\tau)q\sigma_z+(\theta/\tau)\sigma_y$.

\subsection{Four objects}
\label{sec:objects}

\paragraph{A control metric.} The generators of the internal space form the real vector space $\Herm(s)$, which carries the positive Hilbert--Schmidt norm
\begin{equation}
G_{\rm HS}(A)=\frac1s\tr(A^2).
\label{eq:HS}
\end{equation}
Restricted to the span of the coefficients of $H$, it is the natural positive metric on the controls that the symbol couples to translations. Other positive forms of the same type appear in the series: the descent metric $G$ of a gradient controller and the work form $K$ derived from the stress of a field-coupled actuator~\cite{Liu2026QRAIF_B}, and monotone information metrics on states~\cite{Petz1996}.

\paragraph{Causal cones.} The step~\eqref{eq:step} has a strict causal region. After $n$ steps a state supported in $S$ is supported in $S+nN$, the $n$-fold Minkowski sum, so all influence lies in the cone over the polytope $C_N=\tau^{-1}\conv N$. A first-order symbol has a second cone, its characteristic cone, defined in Section~\ref{sec:indefinite}. The two need not agree.

\paragraph{A representation.} The coefficients $\alpha_i,\beta$ generate a real algebra inside $M_s(\mathbb C)$. Whether they satisfy Clifford relations, and with which quadratic form, is a property of the dynamics, not a choice of encoding.

\paragraph{Spectral mass.} For a selected logarithm and energy zero, the positive band's value at $k_*$ is the rest energy $E_0=\hbar|\omega(k_*)|$, and its curvature defines an inertial mass $m_*$ through $\omega(k_*+q)\simeq\omega(k_*)+\hbar q^2/(2m_*)$. A measured phase difference alone does not fix an unknown common energy offset. A Lorentzian read-out identifies $E_0=m_*c^2$, which defines a speed $c$; whether this $c$ agrees with either causal cone is a question about the model.

\begin{proposition}[Positive forms stay positive]
\label{prop:sylvester}
Let $G$ be a positive definite form on a real vector space $V$ and $T:U\to V$ linear. Then $T^*G$, given by $(T^*G)(u,u')=G(Tu,Tu')$, is positive semidefinite, and positive definite if $T$ is injective. In particular no positive definite form is congruent to a form with a negative direction.
\end{proposition}

\begin{proof}
$(T^*G)(u,u)=G(Tu,Tu)\ge0$, with equality only if $Tu=0$. The last statement is Sylvester's law of inertia.
\end{proof}

The Lorentzian form of a read-out is therefore not a pullback or a reparameterization of $G_{\rm HS}$, $G$ or $K$. In the model it arises as a different kind of object, a form on frequency--momentum space built from the spectrum of the evolution.

% ============================================================
\section{Where the indefinite direction comes from}
\label{sec:indefinite}

For a first-order symbol~\eqref{eq:firstorder} write $P(\omega,q)=\det(\omega\mathbf 1-H(q))$ for its characteristic polynomial and $P_0(\omega,q)=\det(\omega\mathbf 1-\sum_iq_i\alpha_i)$ for the principal part. A homogeneous polynomial $p$ on $\mathbb R^n$ is \emph{hyperbolic with respect to} $e$ if $p(e)\neq0$ and $t\mapsto p(\xi+te)$ has only real roots for every $\xi$~\cite{Garding1959}.

\begin{theorem}[Hyperbolicity in the evolution frequency]
\label{thm:hyperbolic}
For Hermitian $\alpha_i,\beta$, every root in $\omega$ of $P(\omega,q)$ is real for every real $q$, and $P_0$ is hyperbolic with respect to the frequency direction $e=(1,0,\dots,0)$. The propagation speed of the principal part in a unit direction $u$ is the largest eigenvalue of $\sum_iu_i\alpha_i$.
\end{theorem}

\begin{proof}
The roots of $P(\cdot,q)$ are the eigenvalues of the Hermitian matrix $H(q)$. For $P_0$, $P_0(e)=1$ and the roots of $t\mapsto P_0(\xi_0+t,\xi')$ are $-\xi_0$ plus the eigenvalues of $\sum_i\xi'_i\alpha_i$, which are real. The characteristic speeds of $i\partial_t\psi=\sum_i\alpha_i(-i\partial_i)\psi$ in direction $u$ are the eigenvalues of $\sum_iu_i\alpha_i$.
\end{proof}

The evolution parameter is supplied by the model. Unitarity makes its conjugate frequency a hyperbolicity direction; it does not choose how many fundamental time parameters a theory should have. The next lemma determines the signature when the characteristic form is quadratic.

\begin{lemma}[Hyperbolic quadratic forms are Lorentzian]
\label{lem:hyperbolic-quadratic}
A real quadratic form $p$ on $\mathbb R^n$ with $p(e)>0$ is hyperbolic with respect to $e$ if and only if it has exactly one positive eigenvalue. Then $p=\ell^2-r$ with $\ell$ linear and $r$ positive semidefinite, and if $p$ is non-degenerate its signature is $(1,n-1)$.
\end{lemma}

\begin{proof}
Decompose $\xi=se+\zeta$ with $\zeta$ in the $p$-orthogonal complement of $e$. Then $p(\xi+te)=p(e)(s+t)^2+p(\zeta)$, which has real roots in $t$ for every $\xi$ if and only if $p(\zeta)\le0$ on the complement. That is the statement that $p$ has one positive direction and is negative semidefinite on a complementary hyperplane.
\end{proof}

\begin{theorem}[Single shell, Clifford relations and the Lorentzian form]
\label{thm:shell}
Let $\alpha_1,\dots,\alpha_d,\beta\in\Herm(s)$ with $\alpha_1,\dots,\alpha_d$ linearly independent, and fix $\mu\in\mathbb R$. The following are equivalent.
\begin{enumerate}
\item[(a)] There is a quadratic form $Q$ on $\mathbb R^d$ such that every eigenvalue $\omega$ of $H(q)$ satisfies $\omega^2=Q(q)+\mu^2$ for all $q$.
\item[(b)] $H(q)^2=(Q(q)+\mu^2)\mathbf 1$ for all $q$.
\item[(c)] $\{\alpha_i,\alpha_j\}=2Q_{ij}\mathbf 1$, $\mu\{\alpha_i,\beta\}=0$ and $\mu^2(\beta^2-\mathbf 1)=0$.
\end{enumerate}
If they hold, then:
\begin{enumerate}
\item $Q$ is positive definite and $Q(q)=G_{\rm HS}\bigl(\sum_iq_i\alpha_i\bigr)$;
\item the read-out form $\eta(\omega,q)=\omega^2-Q(q)$ on frequency--momentum space has signature $(1,d)$, and the characteristic set of $P_0$ lies on its null cone; the sets coincide if $d\ge2$, if $\mu\ne0$, or if $d=1$ and $\alpha_1$ has both positive and negative eigenvalues;
\item if $d+[\mu\neq0]\ge2$, then $s$ is even and $P(\omega,q)=(\omega^2-Q(q)-\mu^2)^{s/2}$;
\item every group velocity $v=\nabla_q\omega$ satisfies $Q^{-1}(v,v)=Q(q)/(Q(q)+\mu^2)\le1$, with equality only for $\mu=0$.
\end{enumerate}
\end{theorem}

\begin{proof}
(a)$\Rightarrow$(b): $H(q)$ is Hermitian, hence diagonalizable, and all its eigenvalues square to $Q(q)+\mu^2$. (b)$\Rightarrow$(c): expanding,
\[
H(q)^2=\sum_{i,j}q_iq_j\tfrac12\{\alpha_i,\alpha_j\}+\mu\sum_iq_i\{\alpha_i,\beta\}+\mu^2\beta^2 .
\]
Comparing the parts of degree two, one and zero in $q$ with $(Q(q)+\mu^2)\mathbf 1$ gives the three relations. (c)$\Rightarrow$(a): the same expansion gives (b), and (b) gives (a).

For item 1, $Q(q)\mathbf 1=(\sum_iq_i\alpha_i)^2\ge0$; if $Q(q)=0$ the Hermitian matrix $\sum_iq_i\alpha_i$ squares to zero and so vanishes, and linear independence gives $q=0$. Taking the normalized trace gives $Q(q)=G_{\rm HS}(\sum_iq_i\alpha_i)$. For item 2, $\eta$ is $\omega^2$ minus a positive definite form, and every root of $P_0(\cdot,q)$ belongs to $\{\pm\sqrt{Q(q)}\}$. When $d\ge2$, the principal symbol at nonzero $q$ anticommutes with an invertible symbol in a $Q$-orthogonal direction, pairing its two signs. If $\mu\ne0$, $\beta$ provides the pairing. In the remaining one-dimensional case both signs occur exactly when $\alpha_1$ has both signs. For item 3, choose coordinates with $Q=c^2\mathbf 1$. The matrices $\alpha_i/c$ together with $\beta$ (if $\mu\neq0$) pairwise anticommute and square to $\mathbf 1$. If at least two are present, each $x$ anticommutes with an invertible $g$, giving $\tr x=\tr(g^{-1}xg)=-\tr x$. Thus $H(q)$ is traceless, and its eigenvalues $\pm\sqrt{Q+\mu^2}$ have equal multiplicity. For item 4, at a differentiable band point with $\omega\ne0$, $\nabla_q\omega=Qq/\omega$ and $Q^{-1}(v,v)=Q(q)/\omega^2$. At a massless node the same statement applies to directional band limits.
\end{proof}

\begin{corollary}[Lorentzian Clifford generators]
\label{cor:gamma}
Under Theorem~\ref{thm:shell} with $\mu\ne0$, the matrices
\begin{equation}
\gamma^0=\beta,\qquad \gamma^i=\beta\alpha_i
\end{equation}
satisfy $\{\gamma^0,\gamma^0\}=2$, $\{\gamma^0,\gamma^i\}=0$ and $\{\gamma^i,\gamma^j\}=-2Q_{ij}$, so they generate a representation of $Cl(1,d)$ on $\mathbb C^s$; $\gamma^0$ is Hermitian and the $\gamma^i$ are anti-Hermitian. The eigenvalue equation $H(q)\psi=\omega\psi$ is equivalent to $(\omega\gamma^0-\sum_iq_i\gamma^i-\mu)\psi=0$, and $(\omega\gamma^0-\sum_iq_i\gamma^i)^2=\eta(\omega,q)\mathbf 1$. For $\mu=0$ and $d\ge1$, the matrices
\begin{equation}
\Gamma^0=\begin{pmatrix}0&\mathbf 1\\ \mathbf 1&0\end{pmatrix},\qquad
\Gamma^i=\begin{pmatrix}0&\alpha_i\\ -\alpha_i&0\end{pmatrix}
\end{equation}
generate a representation of $Cl(1,d)$ on $\mathbb C^s\oplus\mathbb C^s$.
\end{corollary}

\begin{proof}
$(\beta\alpha_i)^2=-\beta^2\alpha_i^2=-Q_{ii}$ and $\{\beta,\beta\alpha_i\}=\alpha_i-\alpha_i=0$ use $\{\alpha_i,\beta\}=0$ and $\beta^2=1$; similarly $\{\beta\alpha_i,\beta\alpha_j\}=-\{\alpha_i,\alpha_j\}$. Multiplying $H(q)-\omega$ by $\beta$ gives the second statement. The block matrices are checked by direct multiplication.
\end{proof}

\paragraph{What has been derived.} In the model the indefinite direction of the read-out form is the frequency $\omega$, the spectral variable of the one-parameter group of time steps. The positive part is the Hilbert--Schmidt norm of the spatial symbol, a control-type metric. The Lorentzian form is the difference $\omega^2-G_{\rm HS}(\sum_iq_i\alpha_i)$, which by Proposition~\ref{prop:sylvester} could not be obtained from any positive form alone. Theorem~\ref{thm:hyperbolic} makes the supplied evolution frequency a hyperbolicity direction for every first-order Hermitian symbol; Lemma~\ref{lem:hyperbolic-quadratic} says that a quadratic cone with one time direction has signature $(1,d)$; and Theorem~\ref{thm:shell} says that the cone is quadratic, with all internal states on it, exactly when the propagation coefficients satisfy Euclidean Clifford relations. The Lorentzian Clifford algebra of Corollary~\ref{cor:gamma} is then constructed, not supplied: the time-like generator is the mass term $\beta$.

The signature is no longer an input. The Clifford relations are equivalent to an observable property, a single shell, which means no birefringence and no internal state that propagates differently from the others. That property, isotropy of $Q$ (which can be absorbed into the choice of spatial coordinates for one species) and the existence of a band point with a linear symbol remain inputs.

\begin{remark}[Two boundary cases]
\label{rem:boundary}
For $d=1$ and $\mu=0$, $\alpha_1^2=Q_{11}\mathbf 1$ allows unequal multiplicities of $\pm Q_{11}^{1/2}$, and for $s=1$ the symbol $H(q)=cq$ has a single branch. Its characteristic set is one chiral branch, with one future null direction. This is the chiral field used as an actuator in Ref.~\cite{Liu2026QRAIF_B}: it supplies one null direction, and a full $(1,1)$ cone needs both chiralities. In the other direction, a quadratic cone can occur without a single shell. The Maxwell field in Riemann--Silberstein form has the symbol $c\,\mathbf S\cdot q$ with spin-one matrices; its characteristic polynomial is $\omega(\omega^2-c^2|q|^2)$, and the flat factor is removed by the Gauss constraint. Theorem~\ref{thm:shell} concerns symbols for which every internal state lies on the cone.
\end{remark}

\subsection{A forward certificate for approximate shells}
\label{sec:residual}

The exact operator identity has a finite-error consequence. It requires bounds on implemented coefficients or their calibration errors, not merely access to effects whose ideal algebra contains those coefficients.

\begin{proposition}[Squared-band certificate]\label{prop:residual}
Choose a positive quadratic form $Q$ and Hermitian coefficients $\alpha_i,\beta$. Let
\[
 R_{ij}=\tfrac12\{\alpha_i,\alpha_j\}-Q_{ij}I,\qquad
 R_{i\beta}=\{\alpha_i,\beta\},\qquad R_\beta=\beta^2-I .
\]
Suppose their operator norms are bounded by calibrated numbers
$\epsilon_{ij},\epsilon_{i\beta},\epsilon_\beta$. If the implemented symbol is
$\widetilde H(q)=H(q)+E(q)$ with Hermitian $E$ and
$\|E(q)\|\le e(q)$, put
\begin{align*}
 \Delta(q)&=\sum_{i,j}|q_iq_j|\epsilon_{ij}
  +|\mu|\sum_i|q_i|\epsilon_{i\beta}+\mu^2\epsilon_\beta,\\
 L(q)&=\sum_i|q_i|\|\alpha_i\|+|\mu|\|\beta\|,\qquad
 B(q)=\Delta(q)+2L(q)e(q)+e(q)^2 .
\end{align*}
Every eigenvalue $\widetilde\omega$ obeys
\[
 \left|\widetilde\omega^2-Q(q)-\mu^2\right|\le B(q).
\]
With $S=Q(q)+\mu^2\ge0$, this gives
$\bigl||\widetilde\omega|-\sqrt S\bigr|\le\sqrt{B(q)}$.
If $S>B(q)$, a sharper radius is
$B(q)/(\sqrt S+\sqrt{S-B(q)})$.
The bounds hold uniformly on any specified window where their input bounds hold.
\end{proposition}
\begin{proof}
Expanding $H^2-SI$ gives $\sum q_iq_jR_{ij}+\mu\sum q_iR_{i\beta}+\mu^2R_\beta$, whose norm is at most $\Delta$. The additional square error is $HE+EH+E^2$, of norm at most $2Le+e^2$. Evaluate on an eigenvector of the Hermitian $\widetilde H$. Finally, $|\sqrt x-\sqrt y|\le\sqrt{|x-y|}$ for nonnegative $x,y$; away from zero use the difference of squares.
\end{proof}

This is a forward dispersion certificate. It does not say that nearly anticommuting coefficients can be repaired to an exact Clifford representation, that the two frequency multiplicities are balanced, or that a finite effect menu implements operator-norm tomography. A calibration valid on a restricted state set is not a norm bound on the entire carrier. The distinction matters: a Hilbert--Schmidt trace identity can hold exactly while different internal states retain different propagation speeds.

\subsection{Outputs to physical control and mass read-out}
\label{sec:physical-interfaces}

Two physical interfaces use the shell theorem, but require different
additional inputs. For control, the spatial map
$J(q)=\sum_iq_i\alpha_i$ is an isometry into internal generators with
$G_{\rm HS}$ because its square is scalar. The trace identity alone
is weaker: the local unitary walk
$W(k)=\operatorname{diag}(e^{-ik},e^{-2ik})$ has
$H(q)=q\operatorname{diag}(1,2)$ and
$G_{\rm HS}(H(q))=5q^2/2$, yet two distinct speeds.
The operator square, not a mean over internal states, supplies one shell.

Q-RAIF B~\cite{Liu2026QRAIF_B} gives that distinction an operational
test. For its centered independent field input, a coupling
$A(a)=J(Ca)/E_*$ has switching work
\[
 W_\rho(a)=\kappa_g\operatorname{tr}[\rho A(a)^2].
\]
If the physical calibration $C$ is invertible, independence of this
charge from every unknown internal state is equivalent to a scalar
spatial shell. Its value is
$\kappa_g Q(Ca)/E_*^2$.
The field profile, calibrated map and energy scale are supplied;
the encounter requires vanishing or commuting free internal
evolution. This is a bridge between two calibrated uses of the same
generators. It neither equates the cost with $\omega^2-Q$ nor
identifies the chiral control field with a general multi-dimensional
propagation carrier. Q-RAIF C~\cite{Liu2026QRAIF_C} uses
Corollary~\ref{cor:gamma}'s symbol to represent the Lorentzian
quadratic space on the closed carrier; discarding a drive gives a
different, reduced-channel representation problem.

For mass--energy read-out, the massive continuum shell must also
carry the positive-energy Poincare representation specified in
Q-RAIF ME~\cite{Liu2026QRAIF_ME}. That paper constructs it on the
positive band of the Dirac symbol, with the boost Jacobian and spin
transport on momentum-space wavefunctions. The symmetric pair of eigenvalues in
Theorem~\ref{thm:shell} alone is not such a representation or a
preparation of its positive sector. A finite lattice window is not
invariant under arbitrary boosts. The actual coin experiment in
Section~\ref{sec:actual-cone} provides a narrower, priced parameter
certificate, which remains useful without upgrading it to exact
relativistic covariance.

% ============================================================
\section{Spatial dimension and the internal space}
\label{sec:dimension}

\begin{theorem}[Dimension bound]
\label{thm:dimension}
Suppose Theorem~\ref{thm:shell} holds on $\mathbb C^s$, and let $n=d+[\mu\neq0]$. Write $s=2^{\nu}m$ with $m$ odd, $\nu=\nu_2(s)$. Then $n\le2\nu+1$. Conversely, for every $s$ and every $n\le2\nu_2(s)+1$ there are coefficients satisfying Theorem~\ref{thm:shell}.
\end{theorem}

\begin{proof}
Choose spatial coordinates in which $Q=c^2\mathbf 1$; the coefficients transform linearly, so the relations are preserved. Then $e_i=\alpha_i/c$ and, if $\mu\ne0$, $e_{d+1}=\beta$ are $n$ pairwise anticommuting Hermitian involutions, which define a representation of the complex Clifford algebra $\mathbb{C}l_n$ on $\mathbb C^s$. This algebra is isomorphic to $M_{2^{n/2}}(\mathbb C)$ for even $n$ and to $M_{2^{(n-1)/2}}(\mathbb C)\oplus M_{2^{(n-1)/2}}(\mathbb C)$ for odd $n$~\cite{LawsonMichelsohn1989}. Every finite-dimensional representation is a direct sum of irreducible ones of dimension $2^{\lfloor n/2\rfloor}$, so $2^{\lfloor n/2\rfloor}$ divides $s$ and $\lfloor n/2\rfloor\le\nu$. For the converse take an irreducible representation of $\mathbb{C}l_{2\nu+1}$, which has dimension $2^\nu$, tensor it with $\mathbf 1_m$, and use any $n$ of its generators.
\end{proof}

\begin{corollary}[Qubit carriers]
\label{cor:qubit}
For $s=2$: an isotropic single-shell symbol exists for massless propagation in $d\le3$ and for massive propagation in $d\le2$. Massive single-shell propagation in $d=3$ requires $4\mid s$; the minimal case $s=4$ is the Dirac symbol $\alpha_i=\sigma_x\otimes\sigma_i$, $\beta=\sigma_z\otimes\mathbf 1$, a pair of Weyl qubits coupled by the mass term.
\end{corollary}

For $s=2$ and $d=3$ the symbol is $H(q)=c\,\sigma\cdot q$ up to a unitary and an orthogonal change of spatial coordinates, and
\begin{equation}
\det\bigl(\omega\mathbf 1-c\,\sigma\cdot q\bigr)=\omega^2-c^2|q|^2 .
\label{eq:weyl}
\end{equation}
This is the Lorentzian form of signature $(1,3)$, obtained from the Pauli relations $\{\sigma_i,\sigma_j\}=2\delta_{ij}$, which are Euclidean. No operator on $\mathbb C^2$ anticommutes with all three Pauli matrices, so $Cl(1,3)$ has no representation on $\mathbb C^2$; the Lorentzian form exists anyway, as the characteristic form of the evolution.

\paragraph{Two forms on the same carrier.}\label{sec:example}
The identity $\det(x_0I+x\cdot\sigma)=x_0^2-|x|^2$ is preserved by
determinant-one filtering, not by the positive Hilbert--Schmidt norm.
For the massive Dirac carrier, the time-like generator in
Corollary~\ref{cor:gamma} is the mass coefficient, not a modular generator.

\begin{proposition}[Full coupling of a qubit]
\label{prop:qubit}
Let $s=2$, $\mu=0$, and let the symbol satisfy Theorem~\ref{thm:shell} with traceless coefficients. The map $q\mapsto\sum_iq_i\alpha_i$ is an injective linear map from $\mathbb R^d$ into the three-dimensional space of traceless Hermitian $2\times2$ matrices. It is onto, so that every internal rotation generator is coupled to some direction of motion, if and only if $d=3$.
\end{proposition}

\begin{proof}
Injectivity is the linear independence assumed in Theorem~\ref{thm:shell}; the target has dimension three.
\end{proof}

Proposition~\ref{prop:qubit} is a precise version of a familiar observation: momentum space of a Weyl particle is identified with the Lie algebra of internal rotations of its qubit. It is a characterization under a stated premise, not a derivation of three dimensions. The premise, that no internal axis is left uncoupled, is a physical input. The following conjecture states what would have to be shown to turn it into a selection principle.

\begin{conjecture}[Minimal carriers and three dimensions]
\label{conj:minimal}
In local unitary dynamics whose elementary internal carrier is a qubit and whose low-energy propagation is isotropic and single-shell, propagation that uses the full internal space selects $d=3$, and massive propagation requires two carriers coupled by a chirality-exchanging term.
\end{conjecture}

Theorem~\ref{thm:dimension} and Proposition~\ref{prop:qubit} reduce this selection proposal to a physical reason for coupling every internal axis to momentum. Irreducibility alone does not supply it: $H=q_x\sigma_x+q_y\sigma_y$ generates the full matrix algebra, while the omitted $\sigma_z$ is not a conserved charge. The three-dimensional Weyl automaton~\cite{DArianoPerinotti2014} and the minimal-directional-system argument of Ref.~\cite{MullerMasanes2013} motivate the proposal, but do not select three dimensions within the present class. A selection principle must distinguish the two-dimensional counterexample by stability, resources or formation, and must explain why a qubit is the relevant elementary carrier.

% ============================================================
\section{The lattice level: causal polytope, three speeds and nodes}
\label{sec:lattice}

\subsection{Group velocities lie in the causal polytope}

\begin{proposition}[Group velocities and the neighbourhood]
\label{prop:polytope}
For the step~\eqref{eq:step}, let $b$ be a band and $k_0$ a point where the eigenvalues of $W(k_0)$ are distinct. Then $\nabla_k\omega_b(k_0)\in\tau^{-1}\conv N$.
\end{proposition}

\begin{proof}
Take $\tau=1$. Let $X$ be the position operator and $\psi$ finitely supported. Then $\langle X\rangle_n=\langle W^n\psi,XW^n\psi\rangle$ satisfies $\langle X\rangle_n-\langle X\rangle_0=\sum_{j<n}\langle W^j\psi,VW^j\psi\rangle$ with $V=W^*[X,W]$. From Eq.~\eqref{eq:step}, $([X,W]\psi)(x)=\sum_hhA_h\psi(x-h)$, so $V$ is bounded and translation invariant with symbol $V(k)=W(k)^*\,i\nabla_kW(k)$. Write $W(k)=\sum_be^{-i\omega_b(k)}P_b(k)$ wherever the eigenvalues are distinct, which is almost everywhere once it holds at one point, since the discriminant of the characteristic polynomial is real-analytic in $k$. For each such $k$ the Ces\`aro means $n^{-1}\sum_{j<n}W(k)^{-j}V(k)W(k)^j$ converge to $\sum_bP_bVP_b$, and differentiating $W=\sum_be^{-i\omega_b}P_b$ with $P_b(\nabla P_b)P_b=0$ gives $P_bVP_b=\nabla\omega_b\,P_b$. By bounded convergence,
\[
\frac{\langle X\rangle_n-\langle X\rangle_0}{n}\longrightarrow\bar v(\psi)=\sum_b\int\nabla\omega_b(k)\,\langle\hat\psi(k),P_b(k)\hat\psi(k)\rangle\,\frac{dk}{|B|},
\]
where $B$ is the Brillouin zone.
On the other hand $W^n\psi$ is supported in $\operatorname{supp}\psi+nN$, so for every unit vector $u$, $u\cdot\langle X\rangle_n\le\max_{x\in\operatorname{supp}\psi}u\cdot x+n\max_{h\in N}u\cdot h$. Hence $u\cdot\bar v(\psi)\le h_N(u):=\max_{h\in N}u\cdot h$. The map $\hat\psi\mapsto\bar v(\psi)$ is continuous in $L^2$ because $|\nabla\omega_b|\le\|V(k)\|$ is bounded, and finitely supported states are dense, so the inequality holds for all $\psi$. Taking $\hat\psi=P_b(k)\varphi$ with $\varphi$ concentrating at $k_0$ gives $u\cdot\nabla\omega_b(k_0)\le h_N(u)$ for every $u$, which is membership in $\conv N$.
\end{proof}

The ballistic scaling used in the proof is the expectation-value form of the weak limit theorem for quantum walks~\cite{GrimmettJansonScudo2004}.

\begin{corollary}[Speed of a massless read-out cone]
\label{cor:inradius}
If a band has a first-order massless symbol at $k_*$ whose cone is round with speed $c$, then $c\le\tau^{-1}r_{\rm in}(\conv N)$, where $r_{\rm in}$ is the inradius.
\end{corollary}

\begin{proof}
Near $k_*$ the group velocity tends to $c\,\hat q$ along every direction $\hat q$, so $c\hat q\in\tau^{-1}\conv N$ for all $\hat q$.
\end{proof}

\subsection{The split-step walk}
\label{sec:walk}

Consider the coined walk that applies a coin rotation and then a conditional shift,
\begin{equation}
W_\theta(k)=\begin{pmatrix}e^{-ika}\cos\theta&-e^{-ika}\sin\theta\\ e^{ika}\sin\theta&e^{ika}\cos\theta\end{pmatrix},\qquad 0<\theta<\frac\pi2,
\label{eq:walk}
\end{equation}
with neighbourhood $N=\{a,-a\}$ and step time $\tau$.

\begin{proposition}[Three speeds of one walk]
\label{prop:walk}
For the walk~\eqref{eq:walk}:
\begin{enumerate}
\item $\cos(\omega\tau)=\cos\theta\cos(ka)$;
\item the support speed is $a/\tau$, and the largest group velocity is $(a/\tau)\cos\theta$, attained at $ka=\pm\pi/2$;
\item near $k=0$, $\omega(k)^2\tau^2=\theta^2+\theta\cot\theta\,(ka)^2+O((ka)^4)$; the rest energy is $E_0=\hbar\theta/\tau$, the inertial mass $m_*=\hbar\tau\tan\theta/a^2$, and the Lorentzian read-out speed defined by $E_0=m_*c_{\rm eff}^2$ is $c_{\rm eff}=(a/\tau)(\theta\cot\theta)^{1/2}$;
\item $\cos\theta<(\theta\cot\theta)^{1/2}<1$.
\end{enumerate}
\end{proposition}

\begin{proof}
Item 1: $\det W_\theta=1$ and $\tr W_\theta=2\cos\theta\cos(ka)$. Item 2: the support statement is Eq.~\eqref{eq:step} with $N=\{\pm a\}$. Differentiating item 1, $v=(a/\tau)\cos\theta\sin(ka)/\sin(\omega\tau)$, and $v^2\le(a/\tau)^2\cos^2\theta$ is equivalent to $\cos^2\theta\cos^2(ka)\le\cos^2(ka)$, with equality only at $\cos(ka)=0$. Item 3: differentiating item 1 twice at $k=0$, where $\omega\tau=\theta$ and $\omega'=0$, gives $\tau\sin\theta\,\omega''=a^2\cos\theta$; then $\tfrac12(\omega^2)''=\omega\omega''$, $m_*=\hbar/\omega''$, and $E_0/m_*=\omega(0)\omega''(0)$. Item 4 is $\sin\theta\cos\theta<\theta<\tan\theta$.
\end{proof}

The Lorentzian speed read from the rest energy and the inertia of slow particles lies strictly between the largest transport speed and the support speed. No wave packet reaches it, because $v\le(a/\tau)\cos\theta<c_{\rm eff}$; and the strict causal cone is wider. In the scaling limit $a,\tau\to0$ with $a/\tau=c$ and $\theta/\tau=\mu$ fixed, all three speeds tend to $c$ and the dispersion tends to $\omega^2=c^2q^2+\mu^2$~\cite{Strauch2006,BisioDArianoTosini2015}. At $\theta=0$ the walk is a pair of exact shifts, with $\omega=\pm k$ and one speed. The deviations at finite mass are $1-c_{\rm eff}^2\tau^2/a^2=\theta^2/3+O(\theta^4)$ and $1-\cos^2\theta=\theta^2+O(\theta^4)$. The consequences for mass read-outs are developed in the companion paper on mass and energy~\cite{Liu2026QRAIF_ME}.

\subsection{The Weyl automaton on the body-centred cubic lattice}
\label{sec:weyl}

The minimal three-dimensional automaton of D'Ariano and Perinotti~\cite{DArianoPerinotti2014} has $s=2$, neighbourhood $N=\{\ell(\pm1,\pm1,\pm1)/\sqrt3\}$ of eight vectors of length $\ell$, and symbol
\begin{equation}
W(k)=e^{-i\kappa_x\sigma_x}\,e^{-i\kappa_y\sigma_y}\,e^{-i\kappa_z\sigma_z},\qquad \kappa=\frac{\ell k}{\sqrt3},
\label{eq:weylqca}
\end{equation}
with $\cos(\omega\tau)=c_xc_yc_z-s_xs_ys_z$, where $c_i=\cos\kappa_i$, $s_i=\sin\kappa_i$. Expanding, $c_xc_yc_z=1-|\kappa|^2/2+O(|\kappa|^4)$ and $s_xs_ys_z=\kappa_x\kappa_y\kappa_z+O(|\kappa|^5)$, so
\begin{equation}
\omega\tau=|\kappa|+\frac{\kappa_x\kappa_y\kappa_z}{|\kappa|}+O(|\kappa|^3).
\label{eq:weylexp}
\end{equation}
At leading order $W(k)\simeq\mathbf 1-i\sigma\cdot\kappa$, the symbol is $c\,\sigma\cdot k$ with $c=\ell/(\sqrt3\tau)$, and Eq.~\eqref{eq:weyl} gives a round Lorentzian cone. Isotropy fails at relative order $|\kappa|$ through the parity-odd term in Eq.~\eqref{eq:weylexp}. Table~\ref{tab:weyl} collects analytic properties and numerical unitary wave-packet diagnostics. It is not a table of records obtained from a physical measuring instrument.

\begin{table}[ht]
\centering
\small
\begin{tabular}{@{}P{0.38\linewidth}P{0.54\linewidth}@{}}
\toprule
\textbf{Property} & \textbf{Value}\\
\midrule
Causal polytope $\tau^{-1}\conv N$ & cube $|x_i|\le\ell/(\sqrt3\tau)$ per unit time (analytic) \\
Inradius of the cube & $\ell/(\sqrt3\tau)=c$: Corollary~\ref{cor:inradius} is saturated (analytic) \\
Group velocity components & at most $c$, attained on the cube faces \\
Largest group speed & $(3/2)^{1/2}c$; the cube corners, at speed $\sqrt3c$, are not reached by group velocities \\
Band-touching points in the Brillouin zone & four: at quasi-energy $0$, $k=0$ (chirality $+1$) and $\kappa=-\tfrac\pi2(1,1,1)$ (chirality $-1$); at quasi-energy $\pi/\tau$, one of each chirality \\
One-site initial state & probability outside the cube: zero to machine precision; outside the round cone of radius $ct$: $0.50$ after $16$ steps, $0.47$ after $24$ \\
Gaussian packets near $k=0$, $24$ steps & probability outside radius $ct+R_0$ ($R_0$ the $99\%$ radius at $t=0$): $0.084$, $0.023$, $0.008$ for momentum widths $\sigma_\kappa=0.4$, $0.25$, $0.15$ \\
\bottomrule
\end{tabular}
\caption{Weyl automaton~\eqref{eq:weylqca}. Chirality is the sign of the determinant of the linear map $q\mapsto$ coefficient of $-i\sigma$ in $\pm W(k_0+q)$. The packets carry a fixed spinor and therefore populate both bands.}
\label{tab:weyl}
\end{table}

\subsection{What the read-out map adds}
\label{sec:readout}

The lattice results show three ways in which the dynamics leaves the Lorentzian read-out underdetermined.

First, at finite mass the speed of the read-out depends on the observable that defines it. In the walk~\eqref{eq:walk}, rest energy and inertia give $c_{\rm eff}$, the fastest wave packets give $(a/\tau)\cos\theta$, and the support front gives $a/\tau$ (Proposition~\ref{prop:walk}).

Second, the Lorentzian cone governs low-momentum states only. In the Weyl automaton a state prepared on one site contains all momenta and leaves almost half of its probability outside the round cone, while it never leaves the cube; the leakage falls as the momentum width shrinks (Table~\ref{tab:weyl}).

Third, the Brillouin zone contains several band-touching points. In the Weyl automaton there are two at quasi-energy zero with opposite chiralities, as the doubling theorem for lattice Hamiltonians would lead one to expect~\cite{NielsenNinomiya1981}, although that theorem does not apply directly to automata. A low-energy observer reads a Weyl cone of one chirality or the other depending on which point the prepared states are near. The dynamics does not choose; the preparation and the instruments do. In one dimension a single chiral cone is possible, since the shift by one site is a legitimate automaton; whether a strictly local isotropic automaton in three dimensions can carry a single node is a question we leave open.

So in this model the read-out map cannot be removed. What can be said precisely is what it consists of: a band point, a momentum window around it, and an observable defining the speed, each with computable consequences. This matches the treatment of access in HAFF, where what a physical reduction resolves depends on the channel and not only on the algebra it generates~\cite{Liu2026PaperA}. Here the Lorentzian geometry read by an observer is the characteristic geometry of the momentum window that its couplings reach.

\subsection{Finite records certify one walk's dictionary}
\label{sec:actual-cone}

The choice of a read-out need not leave all its parameters inaccessible. The companion mass--energy paper realizes the same coin rotation by a fresh Gaussian oscillator, computes its exact open-channel linewidth, and reads its rest phase through controlled mixed-memory instruments~\cite{Liu2026QRAIF_ME}. On a supplied finite periodic ring, a normalized zero mode is preserved by a spatially uniform field coupling. Independent $X$ and $Z$ record trials at one and ten steps resolve the phase alias and certify $\theta$. The ring, uniform support, axes, gain, clock, fresh ports and record archive are priced inputs.

Let that actual protocol give $\theta\in[\theta_-,\theta_+]$ at its stated confidence, with supplied $a,\tau$. Monotonicity of $\tan\theta$, $\cos\theta$ and $C(\theta)=\theta\cot\theta$ gives simultaneous intervals
\begin{align*}
 E_0&\in\frac{\hbar}{\tau}[\theta_-,\theta_+],&
 m_*&\in\frac{\hbar\tau}{a^2}[\tan\theta_-,\tan\theta_+],\\
 c_{\rm eff}^2&\in\frac{a^2}{\tau^2}[C(\theta_+),C(\theta_-)],&
 v_{\max}&\in\frac a\tau[\cos\theta_+,\cos\theta_-].
\end{align*}
At $\theta_+=\pi/2$ the inertial upper bound is infinite. With the priced angle radius $b=.004296823$ at $95\%$ confidence derived in Ref.~\cite{Liu2026QRAIF_ME}, true angle $.3$ and $a=\tau=1$, every confidence interval on that event lies inside $[.3-2b,.3+2b]$. The implied intervals are
\[
 c_{\rm eff}\in[.9838969,.9856636],\qquad
 v_{\max}\in[.9527616,.9578408],\qquad c_{\rm support}=1 .
\]
The support speed is supplied, not measured by the phase protocol. The other two are model implications of its angle certificate. The full preparation, pulse, read-memory and reference budget is given with that protocol.

These are parameter inferences for the prescribed ideal walk, informed by an actual open-drive experiment. They are not its noisy long-time transport velocities: nonzero-momentum transport requires the real Bloch channel and inter-momentum coherences. Nor does calibration prove that the dictionary is Lorentzian throughout the band. The companion finite-angle window theorem gives precisely where a Lorentzian mass reading becomes negative and must be rejected. Thus an observer's interface can be calibrated, can distinguish otherwise conflated speeds, and can fail at a testable boundary.

% ============================================================
\section{Several species: when do the cones become one?}
\label{sec:species}

Theorem~\ref{thm:shell} gives every single-shell species its own form $Q$. For isotropic species in one frame, a common light cone requires their speeds to agree. In a free model nothing makes them agree: each speed is a parameter of the step. This section asks whether coupling the species can make the speeds equal, so that one cone is an output of the dynamics. The answer depends on the kind of coupling. Off-diagonal mixing cannot produce a single quadratic principal shell while its distinct diagonal shells are held fixed. Interactions can change those diagonal propagation laws: in the mediator model their speeds converge when each coupling has divergent accumulated strength toward low energy.

\subsection{Mixing does not merge cones}

Write the principal symbol of two species in a direction $u$ as $A_1(u)$ on $\mathbb C^{s_1}$ and $A_2(u)$ on $\mathbb C^{s_2}$, both single-shell, $A_a(u)^2=Q_a(u)\mathbf 1$. A one-particle coupling adds a block off-diagonal term, which may itself depend linearly on $q$.

\begin{theorem}[Mixing and the single shell]
\label{thm:mixing}
For every $s_1\times s_2$ matrix $V$, the symbol $A=\begin{pmatrix}A_1&V\\ V^*&A_2\end{pmatrix}$ satisfies
\begin{equation}
\min_{\lambda\in\mathbb R}\bigl\|A^2-\lambda\mathbf 1\bigr\|_{\rm HS}^2\ \ge\ \tfrac12\min(s_1,s_2)\,\bigl(Q_1(u)-Q_2(u)\bigr)^2 .
\label{eq:mixing-bound}
\end{equation}
Hence the coupled system has a single shell only if $Q_1=Q_2$. Moreover the largest eigenvalue of $A$ is at least the largest eigenvalue of either block and the smallest is at most the smallest of either block.
\end{theorem}

\begin{proof}
The diagonal blocks of $A^2-\lambda$ are $Q_1-\lambda+VV^*$ and $Q_2-\lambda+V^*V$, and the Hilbert--Schmidt norm of a block matrix is at least that of its diagonal blocks. If $V$ has rank $r$, then $VV^*$ and $V^*V$ share the nonzero eigenvalues $\sigma_1^2,\dots,\sigma_r^2$ and have $s_1-r$ and $s_2-r$ zero eigenvalues. Using $(x+a)^2+(x+b)^2\ge\frac12(a-b)^2$ for each shared eigenvalue, and for $k=\min(s_1,s_2)-r$ pairs of zeros, the squared norm is at least $\frac12(r+k)(Q_1-Q_2)^2$. The last statement is Cauchy interlacing for the principal submatrices $A_1$ and $A_2$.
\end{proof}

The bound~\eqref{eq:mixing-bound} does not depend on the strength of the coupling. A mass-type coupling, independent of $q$, leaves the principal symbol untouched, but can change low-momentum band slopes. For example,
\[
 H(q)=\begin{pmatrix}v_1q&g\\g&v_2q\end{pmatrix},\qquad
 \omega_\pm(q)=\tfrac12(v_1+v_2)q
 \pm\sqrt{g^2+\tfrac14(v_1-v_2)^2q^2}.
\]
For $g\ne0$ both slopes at $q=0$ equal $(v_1+v_2)/2$, although the two distinct principal speeds remain $v_1,v_2$. Equal slopes at one band point are not a common quadratic characteristic cone. An on-site lattice coin can likewise change the effective logarithm of a step; the fixed-block theorem applies only after its principal blocks have been specified. For a chiral principal velocity matrix with fixed diagonal entries, the Schur--Horn theorem makes the eigenvalues majorize those entries, so off-diagonal mixing cannot reduce their variance. The interacting route below instead changes the propagation coefficients themselves.

\subsection{Loops average speeds}

\begin{lemma}[Loops of one speed]
\label{lem:onespeed}
Let a self-energy be a loop whose internal lines all have Lorentzian propagators of speed $c$ and whose vertices are covariant under the Lorentz group of speed $c$. Then the coefficient of $\ln\Lambda$ in it is Lorentz covariant with speed $c$, for any fixed regular cutoff shape enclosing the origin and scaled by $\Lambda$.
\end{lemma}

\begin{proof}
The substitution $k\mapsto(k_0,k/c)$ makes the integrand covariant under the standard Lorentz group. For a homogeneous integrand $f$ of degree $-D$, the coefficient of $\ln\Lambda$ is the flux of the divergence-free field $x f(x)$ through a surface enclosing the origin, and it is the same for every such surface.
\end{proof}

The smallest interacting model is a real scalar of speed $u$ coupled to $N$ Dirac fermions of speed $v$ by $g\phi\bar\psi\psi$ in $3+1$ dimensions. We integrate out a shell of spatial momenta with the frequency unrestricted; the logarithmic coefficients are scheme independent at this order. By Lemma~\ref{lem:onespeed} the fermion loop adds $\frac{Ng^2}{4\pi^2v^3}(p_0^2+v^2p^2)\,dl$ to the inverse scalar propagator $p_0^2+u^2p^2$, where $l=\ln(\Lambda/E)$, so the squared scalar speed is replaced by a convex combination of $u^2$ and $v^2$. The fermion self-energy contains one fermion and one scalar line. Its frequency integrals are
\begin{equation}
J_0=\int\frac{t^2\,dt}{(t^2+v^2)(t^2+u^2)^2}=\frac{\pi}{2u(v+u)^2},\qquad
J_s=\int\frac{dt}{(t^2+v^2)(t^2+u^2)^2}=\frac{\pi(v+2u)}{2vu^3(v+u)^2},
\end{equation}
and the ratio of the spatial to the temporal coefficient gives the speed toward which $v$ moves.

\begin{proposition}[Speed flow in the Yukawa model]
\label{prop:yukawa}
At one loop,
\begin{equation}
\frac{dv}{dl}=a\Bigl(\frac{v+2u}{3}-v\Bigr),\quad a=\frac{g^2}{4\pi^2u(u+v)^2},\qquad
\frac{du^2}{dl}=\frac{Ng^2}{4\pi^2v^3}\,(v^2-u^2),
\label{eq:yukawa-flow}
\end{equation}
and the coupling obeys $dg^2/dl=-(g^4/2\pi^2)\bigl[1/(vu(v+u))+N/(2v^3)+1/(u(v+u)^2)\bigr]$. At $u=v$ the linearized flow of $u-v$ has rate $(6N+1)g^2/(24\pi^2v^3)$, and $u+6Nv$ is conserved to linear order.
\end{proposition}

For $N=1$ these equations coincide with those of Anber and Donoghue~\cite{AnberDonoghue2011}, who derived them by dimensional regularization. The momentum-shell calculation gives the same logarithmic coefficients. Each speed moves toward a target that lies between the speeds of the lines in its loop: the fermion toward $(v+2u)/3$, the scalar toward $v$. The scalar moves $6N$ times faster, so the common speed lies close to the initial fermion speed. Since the Yukawa coupling is marginally irrelevant, $g^2\simeq8\pi^2c_*^3/[(2N+3)l]$, where $c_*$ is the common limiting speed and
\begin{equation}
\frac{|u-v|}{v}\ \propto\ \Bigl[\ln\frac{\Lambda}{E}\Bigr]^{-(6N+1)/(3(2N+3))},
\label{eq:log-rate}
\end{equation}
with exponent $7/15$ for one Dirac fermion. Integrating the full flow to $l=10^9$ gives log--log slopes $-0.4667$ for $N=1$ and $-0.7576$ for $N=4$, as predicted.

\begin{theorem}[Speeds converge through a common mediator]
\label{thm:consensus}
Let fermion species of speeds $v_1,\dots,v_n$ couple to one scalar of speed $u$ with prescribed nonnegative, locally integrable functions $g_1^2(l),\dots,g_n^2(l)$, each fermion obeying the speed equations in~\eqref{eq:yukawa-flow} and $u^2$ receiving the sum of the fermion loops. The theorem is conditional on these coupling histories; a multi-species beta function must separately establish their accumulated strengths. Assume positive initial speeds. Let $D$ be the difference between the largest and the smallest of $u,v_1,\dots,v_n$, and $M_0$ the initial largest speed. Then the largest speed never increases, the smallest never decreases, and
\begin{equation}
D(l)\le D(0)\exp\Bigl(-\frac{1}{24\pi^2M_0^3}\int_0^l\min_ag_a^2\,dl'\Bigr).
\label{eq:consensus}
\end{equation}
If $\int g_a^2\,dl<\infty$, the speed $v_a$ changes in total by at most $D(0)\int g_a^2\,dl/(24\pi^2m_0^3)$, with $m_0$ the initial smallest speed.
\end{theorem}

\begin{proof}
Each $v_a$ moves toward $u$ at rate $\frac23a_a$, and $u^2$ moves toward a weighted mean of the $v_a^2$, so $u$ decreases when it is the largest speed and increases when it is the smallest. Consider the upper Dini derivative $D^+$ of the maximum minus the minimum. If both extremes are fermions $v_a\ge v_b$, then $D^+\le\frac23[a_a(u-v_a)-a_b(u-v_b)]\le-\frac23(\min_ca_c)D$; if the maximum is $u$, then $D^+\le-\frac23a_b(u-v_b)$; the case of $u$ minimal is symmetric. All speeds stay in $[m_0,M_0]$, where $\frac23a_a\ge g_a^2/(24\pi^2M_0^3)$, and Gr\"onwall's inequality gives~\eqref{eq:consensus}. The last statement follows from $|v_a'|\le\frac23a_aD(0)$ and $a_a\le g_a^2/(16\pi^2m_0^3)$.
\end{proof}

The estimate gives a rate when the weakest simultaneous coupling has divergent integral. Convergence itself needs less.

\begin{corollary}[Intermittent interactions]
\label{cor:intermittent}
In the one-loop star model of Theorem~\ref{thm:consensus}, suppose the nonnegative functions $g_a^2$ are locally integrable and initial speeds are positive. Every speed converges. Species connected to the mediator by edges with $\int_0^\infty g_a^2\,dl=\infty$ have the same limiting speed as the mediator. In particular, all speeds converge to one value if each such integral diverges, even when the couplings are never active simultaneously.
\end{corollary}
\begin{proof}
Write $v_a'=A_a(u-v_a)$ and $u'=\sum_aB_a(v_a-u)$, with
\[
 A_a=\frac{g_a^2}{6\pi^2u(u+v_a)^2},\qquad
 B_a=\frac{g_a^2(u+v_a)}{8\pi^2u v_a^3}.
\]
For an active edge, $B_a/A_a=\tfrac34(1+u/v_a)^3$. All speeds stay in $[m_0,M_0]$, so the two directional rates are uniformly comparable, and their integrals diverge exactly when $\int g_a^2$ does. This is a type-symmetric continuous-time consensus system. The cut-balanced convergence theorem~\cite{HendrickxTsitsiklis2013}, applied to the coefficients along the solution, gives convergence and equality of limits within each component of the infinite-integral interaction graph. The theorem allows state-dependent coefficients. It does not provide the same rate as~\eqref{eq:consensus} when $\min_a g_a^2$ vanishes.
\end{proof}

The accumulated interaction across renormalization scales decides the outcome. The variable $l=\ln(\Lambda/E)$ is scale, not laboratory time; intermittent edges here are active on different scale intervals. An identically uncoupled species keeps its cone; a finite integrated coupling need not erase an initial difference. Marginally irrelevant couplings remove differences as a power of $\ln(\Lambda/E)$; near an interacting fixed point they fall as a power of $E/\Lambda$, with the one-loop exponent $7g_*^2/(24\pi^2c_*^3)$ for $N=1$. In $4-\varepsilon$ dimensions the Yukawa fixed point gives $(6N+1)\varepsilon/[3(2N+3)]$ at this order. The weak-coupling model must remain applicable along this flow. A mediator that breaks rotations can instead drive speeds apart, as in nematic criticality of $d$-wave superconductors~\cite{HuhSachdev2008}.

The common speed lies between the smallest and the largest initial speed. Each initial speed of a massless round cone is at most the inradius of $\conv N$ per step (Corollary~\ref{cor:inradius}), so the emergent cone lies inside the causal polytope.

\subsection{Fixed points and the number of cones}

Near an interacting fixed point the question has an answer that does not rely on perturbation theory.

\begin{theorem}[Speed splittings at a conformal fixed point]
\label{thm:cft-cones}
Let the infrared fixed point be a unitary conformal field theory in $D=d+1\ge3$ dimensions, Lorentz invariant for one cone, and consider local perturbations invariant under translations, spatial rotations and CPT. Up to total derivatives the Lorentz-violating perturbations are $\int d^Dx\,O_{0\cdots0}$ for symmetric traceless primaries of even spin $\ell\ge2$. Their couplings satisfy $D-\Delta_O\le2-\ell\le0$; equality in the first inequality means that the current is conserved, while a marginal Lorentz-violating coupling requires a conserved spin-2 current. If the stress tensor is the only conserved spin-2 current, all independent spin-2 speed perturbations are irrelevant. When the excited non-conserved spin-2 modes have a smallest gap $\gamma_2=\min(\Delta_O-D)>0$, their slowest linearized decay is $(E/\Lambda)^{\gamma_2}$.
\end{theorem}

\begin{proof}
Descendants integrate to zero. A primary transforms in a finite-dimensional irreducible representation of the Lorentz group, whose complexification is that of $SO(D)$, and because $(SO(D),SO(D-1))$ is a Gelfand pair only the symmetric traceless representations contain a rotation-invariant component, once each. CPT acts on a Hermitian rank-$\ell$ tensor by $(-1)^\ell$, which removes odd $\ell$, and scalars are Lorentz invariant. The unitarity bound $\Delta\ge\ell+D-2$ for $\ell\ge1$ holds with equality exactly for conserved currents~\cite{Mack1977,Minwalla1998}. The perturbation $\int T_{00}$ rescales time and changes no relative speed.
\end{proof}

The stability argument is established in Refs.~\cite{Sundrum2012,BednikPujolasSibiryakov2013}; it assumes the Lorentz-invariant fixed point and does not prove that every microscopic dynamics reaches it. Additional conserved spin-2 currents are the marginal exceptions. In three dimensions, for exactly two such currents, a suitable basis generates separate translations and the two families of current correlators factorize~\cite[Section 8]{MaldacenaZhiboedov2011}. This is dynamical separation of the current sectors, not necessarily a tensor product of the full theories: a shared singlet constraint can remain. Nor does counting currents count sectors in general; two free scalar fields already have three conserved spin-2 currents. The result needed here is the local criterion for suppressing relative-speed perturbations, while identifying all independently propagating sectors requires further structure.

In two dimensions several propagation sectors can retain different speeds, as in spin--charge separation. A Sugawara decomposition gives an additional sector only when $T-T_J$ is nonzero. Under its positivity, spectrum and finite-signal-speed assumptions, the scale-to-Lorentz theorem permits either one Lorentz factor or several~\cite{Sibiryakov2014}; it does not forbid a common cone from being selected dynamically. Establishing one cone therefore requires information about the actual sectors and their interactions, rather than the dimensionality alone.

\subsection{What the lattice adds}

For number-conserving automata whose interaction preserves the supplied neighbourhood, each distinguishable particle stays in $x_j+nN$ after $n$ steps. The centre-of-mass version of Proposition~\ref{prop:polytope} therefore bounds every isolated bound-state band by $\conv N$ per step. Changing the neighbourhood changes that bound. Equalization of low-energy cones does not itself change the support cone.

For a loop dominated by one linear band point, the logarithmic coefficient carries that point's read-out speed by Lemma~\ref{lem:onespeed}. Finite ultraviolet matching terms need not have the same ratio of temporal and spatial coefficients. A Lorentz-violating regulator can leave corrections of order $g^2/16\pi^2$~\cite{Collins2004}, subsequently reduced only at the infrared rates above.

Composites need not inherit a constituent's read-out cone. Take distinguishable walkers with $\theta_A=.3$, $\theta_B=.5$, $a=\tau=1$, first apply both free coined steps, and then multiply the coincident-position sector by $e^{+.4i}$. With total momentum $K$ conjugate to $(x_A+x_B)/2$, the isolated positive band at $K=0$ has $\omega_0=.7614121$ and $m_*=1/\omega_0''=.9023$, compared with the free threshold $.8$ and constituent mass sum $.8556$. Thus $c_{\rm eff}^2=\omega_0\omega_0''=.8439$, below the constituent values $\theta_A\cot\theta_A=.9698184$ and $\theta_B\cot\theta_B=.9152439$. Periodic relative-coordinate calculations with 91 and 131 sites agree at the quoted precision. They also track a localized band beyond $|K|=.68$; no such endpoint or global maximum speed is inferred from this local mass calculation. Contact interactions can produce bound molecules in quantum walks~\cite{Ahlbrecht2012}. In the example above, the separately computed quasienergy and curvature do not obey the relativistic binding-energy--inertia relation.

A common cone can also be kinematic. The minimal isotropic Weyl automaton is unique on the body-centred cubic lattice~\cite{DArianoPerinotti2014}, so all massless species built from it share $c=\ell/(\sqrt3\tau)$; a Dirac automaton built from two of them has first-order speed $c(1-m^2)^{1/2}$ in lattice units, and photons built from pairs of Weyl automata have the dispersion $2|n_{k/2}|$ and the same low-energy cone~\cite{BisioDArianoPerinotti2016}. For an electron on a lattice of Planck spacing, $m\approx4\times10^{-23}$ and the mass correction is of order $10^{-45}$. Loop matching can generate relative-speed corrections of order $g^2/16\pi^2$ unless a symmetry forbids them. Contact binding also changes the composite cone in the example above, but its size is model dependent and carries no general loop suppression factor.

\subsection{How much running is needed}

With marginal couplings the logarithmic rate~\eqref{eq:log-rate} is far too slow to produce the observed equality of speeds. Reducing a difference by a factor $10^{-13}$ requires $\ln(\Lambda/E)\approx10^{28}$ for one Dirac fermion with $g^2=4\pi$ at the cutoff, and about $10^{12}$ even for a hundred fermions, while $\ln(M_{\rm Pl}/m_e)=51.5$. Anber and Donoghue reached the same qualitative conclusion~\cite{AnberDonoghue2011}.

\begin{conjecture}[One cone needs a strongly coupled regime]
\label{conj:power}
In local unitary dynamics with several interacting gapless species, the relative speed difference of two sectors at scale $E$ is $\delta c(E)\simeq\delta c(\Lambda)\exp(-\int_E^\Lambda\gamma\,d\ln\mu)$, where $\gamma$ is the anomalous dimension of their relative stress tensor and $\delta c(\Lambda)$ is of order $g^2/16\pi^2$ from lattice matching even when the step is common. Equalities of order $10^{-15}$ then require $\int\gamma\,d\ln\mu\gtrsim30$ above accessible energies, that is, a strongly coupled regime with $\gamma$ of order one over many decades, or a symmetry that removes $\delta c(\Lambda)$. The largest residual differences sit between the most weakly coupled sectors.
\end{conjecture}

Theorems~\ref{thm:consensus} and~\ref{thm:cft-cones} tie equalization to interaction strength; lattice matching regenerates relative-speed differences, while marginal running removes them only logarithmically. The quantitative estimates above and the power-law emergence in strongly coupled holographic models~\cite{BednikPujolasSibiryakov2013} support the conjecture. Anber and Donoghue's observation that the weakest interactions leave the largest differences also makes gravitational waves against light a sensitive comparison~\cite{AnberDonoghue2011}.

A viable ultraviolet realization must place the strongly coupled range consistently with the observed weakly coupled particle description and experimental bounds. A composite realization is one possibility, not a necessary consequence of the conjecture. A symmetry route, such as non-relativistic supersymmetry protecting dimension-four terms, can avoid strong coupling but must survive symmetry breaking. Gravity is shared by all species yet is irrelevant at low energy; this flow cannot align its cone there. The relation of the gravitational-wave and matter cones therefore requires an ultraviolet mechanism or a protecting symmetry.

\subsection{Consequences}

The mediator model makes a common speed a conditional output of divergent accumulated interactions. The conformal argument establishes stability against relative-speed perturbations near a supplied Lorentz-invariant fixed point; it does not establish formation of that fixed point. This separates three inputs to the higher-dimensional Consciousness construction: each species' shell and signature, the alignment of their cones, and the actual infrared regime. Its boosted-ball family requires a common cone. In the two-dimensional HAFF gravity model~\cite{Liu2026HAFFD}, identifying the matter and Jackiw--Teitelboim cones remains an input; the two-dimensional sector theorem neither supplies nor forbids that identification.

% ============================================================
\section{Relations to other frameworks and papers}
\label{sec:relations}

\subsection{Algebraic frameworks}

Sequential operations are represented here by linear maps under ordinary composition, which is associative by definition. Jordan and octonionic models need additional bracketing and composition rules before tensor products, subsystems or error correction can be transferred~\cite{Schafer1966,Gunaydin1973,AlfsenShultz2003}; this is a statement of scope, not a no-go theorem. Von Neumann and $C^*$-algebras are compatible with everything above, since the step~\eqref{eq:step} and its symbol live in a concrete $C^*$-algebra; such algebras do not single out a quadratic form until a dynamics is given. Clifford algebras are distinguished in this paper by a dynamical property, the single shell of Theorem~\ref{thm:shell}, rather than by being chosen as an encoding.

\subsection{Thermodynamic gravity, modular flow and spectral triples}

Jacobson's derivation of the Einstein equation assumes local Lorentzian causal structure, local equilibrium and an entropy--area relation~\cite{Jacobson1995}, and entanglement derivations of linearized gravity work on a Lorentzian background~\cite{Faulkner2014,Lashkari2014}. The present model supplies local Lorentzian cones for the low-momentum sector of a local unitary dynamics. It supplies neither equilibrium nor the entropy--area relation, and Section~\ref{sec:readout} shows that its cones depend on the momentum window.

Local algebras in relativistic quantum field theory are of Type~III~\cite{Haag1996,Witten2018}. For the vacuum and wedge regions the Bisognano--Wichmann theorem identifies modular flow with the boost subgroup~\cite{BisognanoWichmann1975,BisognanoWichmann1976}; the identification presupposes the Lorentzian structure it describes. Whether the modular data of states of a lattice automaton reproduce the boosts of its low-energy cone in the scaling limit is open and could be computed in the walk of Section~\ref{sec:walk}.

In Connes' noncommutative geometry a Dirac operator with Euclidean Clifford symbol encodes Riemannian geometry~\cite{Connes1994}, and Lorentzian versions are an active problem~\cite{Barrett2007,PaschkeSitarz2006,VandenDungen2016}. Theorem~\ref{thm:shell} says that using such a Dirac operator as a Hamiltonian, $i\partial_t\psi=D\psi$, produces the Lorentzian characteristic cone of signature $(1,d)$ on time times space. We have not worked out the relation to the Krein-space approaches cited.

\subsection{The Q-RAIF papers and HAFF}

Section~\ref{sec:physical-interfaces} specifies the separate outputs to
control and spectral mass. The actuator field of Q-RAIF B is the
single-chirality case of Remark~\ref{rem:boundary}; its outgoing stress
is used in HAFF-D~\cite{Liu2026HAFFD}. Identifying that matter cone with
the gravitational cone is a further physical input, not a consequence
of an isometry of internal generators.

% ============================================================
\section{Discussion}
\label{sec:discussion}

\begin{table}[ht]
\centering
\small
\begin{tabular}{@{}P{0.28\linewidth}P{0.62\linewidth}@{}}
\toprule
\textbf{Ingredient of a Lorentzian read-out} & \textbf{Status in local unitary propagation}\\
\midrule
Evolution time & supplied; unitarity makes its frequency a hyperbolicity direction (Theorem~\ref{thm:hyperbolic}) \\
Signature $(1,d)$ & derived for any quadratic cone (Lemma~\ref{lem:hyperbolic-quadratic}) \\
Clifford relations & equivalent to a single shell (Theorem~\ref{thm:shell}); Lorentzian generators constructed (Corollary~\ref{cor:gamma}) \\
Single shell, isotropy, linear band point & inputs; observable and testable in a given model \\
Spatial dimension $d$ & input; bounded by the internal space (Theorem~\ref{thm:dimension}) \\
Speed $c$ & a property of the dynamics; at finite mass it depends on the defining observable (Proposition~\ref{prop:walk}) \\
Read-out map & input: band point, momentum window and speed observable (Section~\ref{sec:readout}) \\
Common speed of several species & conditional output of the mediator model; stable against irrelevant relative-speed perturbations near a supplied Lorentz-invariant fixed point, whose formation is not derived (Theorems~\ref{thm:consensus}, \ref{thm:cft-cones}) \\
Lattice, neighbourhood, step & inputs \\
\bottomrule
\end{tabular}
\caption{Inputs of a Lorentzian read-out in local unitary propagation.}
\label{tab:status}
\end{table}

Table~\ref{tab:status} separates the derived signature from the remaining inputs: lattice data, absence of birefringence and the observer's read-out map. Deriving the first two requires a formation model; no property of the dynamics alone specifies the third.

The problems that follow have the most leverage. Section~\ref{sec:species} answers the question of a common cone for several species under stated conditions; what remains there is the anomalous dimension $\gamma_2$ in a strongly coupled lattice model, which Conjecture~\ref{conj:power} needs. Another is chirality: whether a strictly local, isotropic three-dimensional automaton can have a single node, which would remove one component of the read-out map. The mass-setting interface now has a nonempty field realization and finite certificate, but renewed drives, full noisy spatial transport and covariant support remain physical tasks. A fourth is the modular route: computing the modular flow of a lattice state and comparing it with the boosts of the low-energy cone. Position-dependent coins, which produce effective curved metrics, would connect the model to the gravitational questions of HAFF.

\section{Conclusion}

Propagation supplies the relative sign in $\omega^2-G_{\rm HS}$; a positive control metric alone cannot. The single-shell criterion, dimension bound and calibrated residuals connect this form to a specified carrier. Finite angle records imply separated ideal dictionary-speed intervals, while noisy transport remains a different observable. For several species, the mediator model aligns cones under accumulated-coupling conditions; the conformal argument tests their stability, not the creation of a fixed point. These are the respective interfaces to mass--energy read-out, the common-cone geometry of Consciousness II and the still-supplied matter cone of HAFF-D. The calibrated work bridge constrains how propagation generators can be reused for control, without turning positivity into an indefinite signature. Formation of the carrier, clock and accessible momentum window remains a dynamical question.

% ============================================================
\begin{thebibliography}{99}
\raggedright

\bibitem{HendrickxTsitsiklis2013}
J.~M.~Hendrickx and J.~N.~Tsitsiklis, \emph{Convergence of Type-Symmetric and Cut-Balanced Consensus Seeking Systems}, IEEE Trans.\ Autom.\ Control \textbf{58}, 214--218 (2013), arXiv:1102.2361.

\bibitem{RyuTakayanagi2006}
S.~Ryu and T.~Takayanagi, \emph{Holographic Derivation of Entanglement Entropy from AdS/CFT}, Phys.\ Rev.\ Lett.\ \textbf{96}, 181602 (2006).

\bibitem{Jacobson1995}
T.~Jacobson, \emph{Thermodynamics of Spacetime: The Einstein Equation of State}, Phys.\ Rev.\ Lett.\ \textbf{75}, 1260 (1995).

\bibitem{Hestenes1966}
D.~Hestenes, \emph{Space-Time Algebra}, Gordon and Breach (1966).

\bibitem{Doran2003}
C.~Doran and A.~Lasenby, \emph{Geometric Algebra for Physicists}, Cambridge University Press (2003).

\bibitem{Dirac1928}
P.~A.~M.~Dirac, \emph{The quantum theory of the electron}, Proc.\ R.\ Soc.\ Lond.\ A \textbf{117}, 610--624 (1928).

\bibitem{Garding1959}
L.~G{\aa}rding, \emph{An inequality for hyperbolic polynomials}, J.\ Math.\ Mech.\ \textbf{8}, 957--965 (1959).

\bibitem{LawsonMichelsohn1989}
H.~B.~Lawson and M.-L.~Michelsohn, \emph{Spin Geometry}, Princeton University Press (1989).

\bibitem{BialynickiBirula1994}
I.~Bia{\l}ynicki-Birula, \emph{Weyl, Dirac, and Maxwell equations on a lattice as unitary cellular automata}, Phys.\ Rev.\ D \textbf{49}, 6920--6927 (1994).

\bibitem{Meyer1996}
D.~A.~Meyer, \emph{From quantum cellular automata to quantum lattice gases}, J.\ Stat.\ Phys.\ \textbf{85}, 551--574 (1996).

\bibitem{Strauch2006}
F.~W.~Strauch, \emph{Relativistic quantum walks}, Phys.\ Rev.\ A \textbf{73}, 054302 (2006).

\bibitem{BisioDArianoTosini2015}
A.~Bisio, G.~M.~D'Ariano, and A.~Tosini, \emph{Quantum field as a quantum cellular automaton: The Dirac free evolution in one dimension}, Ann.\ Phys.\ \textbf{354}, 244--264 (2015).

\bibitem{DArianoPerinotti2014}
G.~M.~D'Ariano and P.~Perinotti, \emph{Derivation of the Dirac equation from principles of information processing}, Phys.\ Rev.\ A \textbf{90}, 062106 (2014), arXiv:1306.1934.

\bibitem{RaetzelRiveraSchuller2011}
D.~R\"atzel, S.~Rivera, and F.~P.~Schuller, \emph{Geometry of physical dispersion relations}, Phys.\ Rev.\ D \textbf{83}, 044047 (2011), arXiv:1010.1369.

\bibitem{MullerMasanes2013}
M.~P.~M\"uller and L.~Masanes, \emph{Three-dimensionality of space and the quantum bit: an information-theoretic approach}, New J.\ Phys.\ \textbf{15}, 053040 (2013), arXiv:1206.0630.

\bibitem{GrimmettJansonScudo2004}
G.~Grimmett, S.~Janson, and P.~F.~Scudo, \emph{Weak limits for quantum random walks}, Phys.\ Rev.\ E \textbf{69}, 026119 (2004).

\bibitem{NielsenNinomiya1981}
H.~B.~Nielsen and M.~Ninomiya, \emph{Absence of neutrinos on a lattice (I). Proof by homotopy theory}, Nucl.\ Phys.\ B \textbf{185}, 20--40 (1981).

\bibitem{AnberDonoghue2011}
M.~M.~Anber and J.~F.~Donoghue, The emergence of a universal limiting speed, \emph{Phys.\ Rev.\ D} \textbf{83}, 105027 (2011), arXiv:1102.0789.

\bibitem{Ahlbrecht2012}
A.~Ahlbrecht, A.~Alberti, D.~Meschede, V.~B.~Scholz, A.~H.~Werner and R.~F.~Werner,
\emph{Molecular binding in interacting quantum walks}, New J.\ Phys.\ \textbf{14}, 073050 (2012).

\bibitem{HuhSachdev2008}
Y.~Huh and S.~Sachdev, Renormalization group theory of nematic ordering in $d$-wave superconductors, \emph{Phys.\ Rev.\ B} \textbf{78}, 064512 (2008), arXiv:0806.0002.

\bibitem{Mack1977}
G.~Mack, All unitary ray representations of the conformal group SU(2,2) with positive energy, \emph{Commun.\ Math.\ Phys.} \textbf{55}, 1--28 (1977).

\bibitem{Minwalla1998}
S.~Minwalla, Restrictions imposed by superconformal invariance on quantum field theories, \emph{Adv.\ Theor.\ Math.\ Phys.} \textbf{2}, 781--846 (1998), arXiv:hep-th/9712074.

\bibitem{Sundrum2012}
R.~Sundrum, From fixed points to the fifth dimension, \emph{Phys.\ Rev.\ D} \textbf{86}, 085025 (2012), arXiv:1106.4501.

\bibitem{BednikPujolasSibiryakov2013}
G.~Bednik, O.~Pujol\`as and S.~Sibiryakov, Emergent Lorentz invariance from strong dynamics: holographic examples, \emph{JHEP} \textbf{11} (2013) 064, arXiv:1305.0011.

\bibitem{MaldacenaZhiboedov2011}
J.~Maldacena and A.~Zhiboedov, Constraining conformal field theories with a higher spin symmetry, \emph{J.\ Phys.\ A} (2013), arXiv:1112.1016.

\bibitem{Sibiryakov2014}
S.~Sibiryakov, From scale invariance to Lorentz symmetry, \emph{Phys.\ Rev.\ Lett.} \textbf{112}, 241602 (2014), arXiv:1403.4742.

\bibitem{BisioDArianoPerinotti2016}
A.~Bisio, G.~M.~D'Ariano and P.~Perinotti, Quantum cellular automaton theory of light, \emph{Ann.\ Phys.} \textbf{368}, 177--190 (2016), arXiv:1407.6928.

\bibitem{Collins2004}
J.~Collins, A.~Perez, D.~Sudarsky, L.~Urrutia, and H.~Vucetich, \emph{Lorentz invariance and quantum gravity: an additional fine-tuning problem?}, Phys.\ Rev.\ Lett.\ \textbf{93}, 191301 (2004).

\bibitem{Petz1996}
D.~Petz, \emph{Monotone metrics on matrix spaces}, Linear Algebra Appl.\ \textbf{244}, 81--96 (1996).

\bibitem{Schafer1966}
R.~D.~Schafer, \emph{An Introduction to Nonassociative Algebras}, Academic Press (1966).

\bibitem{Gunaydin1973}
M.~G\"unaydin and F.~G\"ursey, \emph{Quark structure and octonions}, J.\ Math.\ Phys.\ \textbf{14}, 1651 (1973).

\bibitem{AlfsenShultz2003}
E.~M.~Alfsen and F.~W.~Shultz, \emph{Geometry of State Spaces of Operator Algebras}, Birkh\"auser (2003).

\bibitem{Faulkner2014}
T.~Faulkner, M.~Guica, T.~Hartman, R.~C.~Myers, and M.~Van~Raamsdonk, \emph{Gravitation from Entanglement in Holographic CFTs}, JHEP \textbf{03}, 051 (2014).

\bibitem{Lashkari2014}
N.~Lashkari, M.~B.~McDermott, and M.~Van~Raamsdonk, \emph{Gravitational dynamics from entanglement ``thermodynamics''}, JHEP \textbf{04}, 195 (2014).

\bibitem{Haag1996}
R.~Haag, \emph{Local Quantum Physics: Fields, Particles, Algebras}, Springer-Verlag (1996).

\bibitem{Witten2018}
E.~Witten, \emph{APS Medal for Exceptional Achievement in Research: Invited article on entanglement properties of quantum field theory}, Rev.\ Mod.\ Phys.\ \textbf{90}, 045003 (2018).

\bibitem{BisognanoWichmann1975}
J.~J.~Bisognano and E.~H.~Wichmann, \emph{On the duality condition for a Hermitian scalar field}, J.\ Math.\ Phys.\ \textbf{16}, 985--1007 (1975).

\bibitem{Connes1994}
A.~Connes, \emph{Noncommutative Geometry}, Academic Press (1994).

\bibitem{Barrett2007}
J.~W.~Barrett, \emph{A Lorentzian version of the non-commutative geometry of the standard model of particle physics}, J.\ Math.\ Phys.\ \textbf{48}, 012303 (2007).

\bibitem{PaschkeSitarz2006}
M.~Paschke and A.~Sitarz, \emph{Equivariant Lorentzian spectral triples}, arXiv:math-ph/0611029 (2006).

\bibitem{VandenDungen2016}
K.~van den Dungen, \emph{Krein spectral triples and the fermionic action}, Math.\ Phys.\ Anal.\ Geom.\ \textbf{19}, 4 (2016).

\bibitem{Liu2026PaperA}
S.~Liu, \emph{Physical Coarse-Graining and the Problem of Emergent Geometry: The Holographic Alaya--Field Framework}, (2026), series record DOI: 10.5281/zenodo.18361706.

\bibitem{Liu2026PaperB}
S.~Liu, \emph{Coupling, Records, and Stable Relations: Formation, Protected Content and the Lifetime of a Geometry}, (2026), series record DOI: 10.5281/zenodo.18367060.

\bibitem{Liu2026HAFFD}
S.~Liu, \emph{Selection, Maintenance and Gravitational Response: Modular Constraints and a Physical Return}, (2026), series record DOI: 10.5281/zenodo.18388881.

\bibitem{Liu2026QRAIF_B}
S.~Liu, \emph{Thermodynamic Stability Constraints on Persistent Open Quantum Subsystems: Control Algebra, Physical Actuation and Energy Cost}, (2026), series record DOI: 10.5281/zenodo.18525890.

\bibitem{Liu2026QRAIF_C}
S.~Liu, \emph{Representing Control in the World: Retained History, Physical Sources and Renewal in Q-RAIF}, (2026), series record DOI: 10.5281/zenodo.18528934.

\bibitem{Liu2026QRAIF_ME}
S.~Liu, \emph{Recovering Mass--Energy Equivalence within Algebraic Lorentzian Read-out: Spectral Mass, Lattice Windows and Finite Records}, (2026), series record DOI: 10.5281/zenodo.20280951.

\bibitem{BisognanoWichmann1976}
J.~J.~Bisognano and E.~H.~Wichmann, \emph{On the duality condition for quantum fields}, J.\ Math.\ Phys.\ \textbf{17}, 303--321 (1976), doi:10.1063/1.522898.

\end{thebibliography}

\end{document}
